\documentclass[11pt,a4paper]{article}
\usepackage[utf8]{inputenc}
\usepackage[T1]{fontenc}
\usepackage{amsmath, amssymb}
\usepackage{amsthm}
\usepackage{booktabs}
\usepackage[margin=2.5cm]{geometry}
\usepackage{hyperref}
\hypersetup{pdftitle={Energy Decay in a Damped Harmonic Oscillator}, pdfauthor={Lena Ortiz}}

\newtheorem{theorem}{Theorem}

\title{Energy Decay in a Damped Harmonic Oscillator}
\author{Lena Ortiz\\ \small Department of Physics, Northfield College}
\date{October 2026}

\begin{document}
\maketitle

\begin{abstract}
We study the free motion of a mass on a linear spring subject to viscous damping. We derive the closed-form solution in the underdamped regime, show that the mechanical energy decays monotonically, and compare the analytic solution with two simple numerical schemes.
\end{abstract}

\section{Introduction}\label{sec:intro}

The damped harmonic oscillator is the simplest model of a system that loses energy to its surroundings~\cite{ref1}. Its equation of motion is
\begin{equation}
  m\ddot{x}+c\dot{x}+kx=0,
  \label{eq:gauss}
\end{equation}
where $m>0$ is the mass, $c\geq0$ the damping coefficient and $k>0$ the spring constant.\footnote{All quantities are in SI units.} Dividing by $m$ and writing $\omega_{0}=\sqrt{k/m}$ and $\gamma=c/(2m)$ gives the \emph{standard form}
\begin{equation}
  \ddot{x}+2\gamma\dot{x}+\omega_{0}^{2}x=0.
  \label{eq:2}
\end{equation}

\section{Solution}

Substituting $x=e^{\lambda t}$ into~\eqref{eq:2} gives the characteristic equation $\lambda^{2}+2\gamma\lambda+\omega_{0}^{2}=0$, whose roots are
\begin{equation}
  \lambda_{\pm}=-\gamma\pm\sqrt{\gamma^{2}-\omega_{0}^{2}}.
\end{equation}
The behaviour depends on the sign of the discriminant:
\begin{equation*}
  \text{regime}=\begin{cases}
\text{underdamped,} & \gamma<\omega_{0}, \\
\text{critically damped,} & \gamma=\omega_{0}, \\
\text{overdamped,} & \gamma>\omega_{0}.
\end{cases}
\end{equation*}
In the underdamped case let $\omega_{d}=\sqrt{\omega_{0}^{2}-\gamma^{2}}$. With $x(0)=x_{0}$ and $\dot{x}(0)=0$ the solution is
\begin{equation}
  x(t)=x_{0}e^{-\gamma t}\left( \cos\omega_{d}t+\frac{\gamma}{\omega_{d}}\sin\omega_{d}t \right).
  \label{eq:cos-dsin}
\end{equation}
Equivalently, with $\mathbf{y}=(x,\dot{x})$, equation~\eqref{eq:2} becomes the first-order system
\begin{equation}
  \dot{\mathbf{y}}=\begin{pmatrix}
0 & 1 \\
-\omega_{0}^{2} & -2\gamma
\end{pmatrix}\mathbf{y}.
\end{equation}

\begin{theorem}
  The energy $E(t)=\frac{1}{2}m\dot{x}^{2}+\frac{1}{2}kx^{2}$ never increases. More precisely, $\frac{dE}{dt}=-c\dot{x}^{2}\leq0$.
\end{theorem}

\begin{proof}
  Differentiate $E$ and use~\eqref{eq:gauss}:
\end{proof}

\begin{align*}
  \frac{dE}{dt} &=m\dot{x}\ddot{x}+kx\dot{x} \\
   &=\dot{x}(m\ddot{x}+kx)=-c\dot{x}^{2}.
\end{align*}

\section{Numerical check}

We integrated~\eqref{eq:2} with $\omega_{0}=2$, $\gamma=0.2$ and $x_{0}=1$ on $t\in[0,10]$, using the explicit Euler method and the classical fourth-order Runge--Kutta method (RK4)~\cite{ref2}.

For each step size $h$ we:
\begin{enumerate}
  \item integrated from $t=0$ to $t=10$ with $10/h$ steps;
  \item compared every step with the exact solution~\eqref{eq:cos-dsin};
  \item recorded the largest absolute difference.
\end{enumerate}
Table~\ref{tab:maximum-absolute} lists the results and Figure~\ref{fig:the-solution} shows the exact solution.

\begin{table}[htbp]
  \centering
  \caption{Maximum absolute error on $[0,10]$.}
  \label{tab:maximum-absolute}
  \begin{tabular}{lcc}
    \toprule
    \textbf{Step $h$} & \textbf{Euler} & \textbf{RK4} \\
    \midrule
    0.100 & $8.51\times10^{-1}$ & $5.00\times10^{-5}$ \\
    0.050 & $2.50\times10^{-1}$ & $3.10\times10^{-6}$ \\
    0.025 & $1.06\times10^{-1}$ & $1.93\times10^{-7}$ \\
    \bottomrule
  \end{tabular}
\end{table}

\begin{figure}[htbp]
  \centering
  \fbox{\parbox[c][0.42\linewidth][c]{0.7\linewidth}{\centering\small\itshape Image not chosen yet}}
  \caption{The solution~\eqref{eq:cos-dsin} for $\gamma=0.2$, $\omega_{0}=2$ and its envelope $\pm e^{-\gamma t}$.}
  \label{fig:the-solution}
\end{figure}

Two things stand out:
\begin{itemize}
  \item Euler's error is large because the method \textbf{adds} energy at every step, which works against the physical damping.
  \item Halving $h$ divides the RK4 error by about $16=2^{4}$, as expected for a fourth-order method.
\end{itemize}

\begin{thebibliography}{9}
\bibitem{ref1} H. Goldstein, C. Poole and J. Safko, \emph{Classical Mechanics}, 3rd ed., Addison-Wesley, 2002.
\bibitem{ref2} E. Hairer, S. P. Nørsett and G. Wanner, \emph{Solving Ordinary Differential Equations I}, 2nd ed., Springer, 1993.
\end{thebibliography}

\end{document}
